% GENERATED FILE. Edit canonical lessons, then run npm run generate:books.
% canonical-source-hash: fnv1a64:b6e91dc827a6c924
% canonical-lessons: ES-C35-sintesis-gustar

\chapter{Synthesis --- What You Like, and Who Does the Liking}
\label{ch:spanish-synthesis-gustar}
\hlchaptermodality{80}

\begin{chapteropening}
\textbf{By the end of this chapter:} \emph{I can say what I like without translating I like word by word, and hear mucho gusto as the same word.}

Liking started at the tongue; and the chapter-5 phrase turns out to be gustar's noun.
\end{chapteropening}

\section[Practice]{Synthesis — what you like, and who is doing the liking}

\label{lesson:ES-C35-sintesis-gustar}



\begin{quote}
Say \textquotedblleft{}I like coffee\textquotedblright{} in Spanish. Then answer this: in that sentence, which word is the \textbf{subject}?
\end{quote}

\begin{grammarlens}[title={Grammar Lens: the sentence that will not let you be the subject}]
\textbf{Me gusta el café.} The coffee is the subject. The coffee does the pleasing. You are the one it happens \emph{to}.

English puts you in charge: \emph{I like coffee} — I am the subject, the coffee just sits there. Spanish reverses it.

\begin{quote}
Two chapters ago, all three verbs were about \textbf{the one doing it}. That is why this chapter had to sit next to them.
\end{quote}

Translate \emph{I like} word by word and you build a sentence Spanish does not have — which is why this is a chapter about saying an idea, not a word.
\end{grammarlens}

\subsection*{The exchange}
\begin{quote}
— ¿Te gusta el café? — Sí, \textbf{me gusta} mucho. ¿Y a ti? — A mí también. \textbf{Mucho gusto.}
\end{quote}

That last line is a joke you can now get. \emph{Mucho gusto} — the phrase you learned long ago for \textquotedblleft{}pleased to meet you\textquotedblright{} — is the \textbf{same word}. It has always meant \emph{much pleasure}. You have been saying \emph{gustar}'s noun since your third conversation in this book.

\begin{grammarlens}[title={What you've built}]
\emph{Gustar} is Latin \emph{gustāre}, \textbf{to taste} — English \emph{gusto}, \emph{disgust}. Liking, in Spanish, started at the tongue.
\end{grammarlens}

\subsection*{Your turn}
Say these aloud:

\begin{itemize}
\raggedright
  \item the exchange above, both parts
  \item it again, and as you say \emph{me gusta}, notice that you are \textbf{not} the subject
\end{itemize}

\emph{Twice through:} Until \emph{me gusta} arrives without translating \emph{I like} first.

\begin{tcolorbox}[breakable,colback=teal!4,colframe=teal!35!black,title={Before you move on}]
Who is the subject of \emph{me gusta el café}? (\textbf{El café}.) What did \emph{gustar} mean first? (\textbf{To taste}.) And what has \emph{mucho gusto} meant all along? (\textbf{Much pleasure} — the same word.)
\end{tcolorbox}
